% Options for packages loaded elsewhere
\PassOptionsToPackage{unicode}{hyperref}
\PassOptionsToPackage{hyphens}{url}
\documentclass[
]{article}
\usepackage{xcolor}
\usepackage{amsmath,amssymb}
\setcounter{secnumdepth}{5}
\usepackage{iftex}
\ifPDFTeX
  \usepackage[T1]{fontenc}
  \usepackage[utf8]{inputenc}
  \usepackage{textcomp} % provide euro and other symbols
\else % if luatex or xetex
  \usepackage{unicode-math} % this also loads fontspec
  \defaultfontfeatures{Scale=MatchLowercase}
  \defaultfontfeatures[\rmfamily]{Ligatures=TeX,Scale=1}
\fi
\usepackage{lmodern}
\ifPDFTeX\else
  % xetex/luatex font selection
\fi
% Use upquote if available, for straight quotes in verbatim environments
\IfFileExists{upquote.sty}{\usepackage{upquote}}{}
\IfFileExists{microtype.sty}{% use microtype if available
  \usepackage[]{microtype}
  \UseMicrotypeSet[protrusion]{basicmath} % disable protrusion for tt fonts
}{}
\makeatletter
\@ifundefined{KOMAClassName}{% if non-KOMA class
  \IfFileExists{parskip.sty}{%
    \usepackage{parskip}
  }{% else
    \setlength{\parindent}{0pt}
    \setlength{\parskip}{6pt plus 2pt minus 1pt}}
}{% if KOMA class
  \KOMAoptions{parskip=half}}
\makeatother
\usepackage{longtable,booktabs,array}
\usepackage{calc} % for calculating minipage widths
% Correct order of tables after \paragraph or \subparagraph
\usepackage{etoolbox}
\makeatletter
\patchcmd\longtable{\par}{\if@noskipsec\mbox{}\fi\par}{}{}
\makeatother
% Allow footnotes in longtable head/foot
\IfFileExists{footnotehyper.sty}{\usepackage{footnotehyper}}{\usepackage{footnote}}
\makesavenoteenv{longtable}
\usepackage{graphicx}
\makeatletter
\newsavebox\pandoc@box
\newcommand*\pandocbounded[1]{% scales image to fit in text height/width
  \sbox\pandoc@box{#1}%
  \Gscale@div\@tempa{\textheight}{\dimexpr\ht\pandoc@box+\dp\pandoc@box\relax}%
  \Gscale@div\@tempb{\linewidth}{\wd\pandoc@box}%
  \ifdim\@tempb\p@<\@tempa\p@\let\@tempa\@tempb\fi% select the smaller of both
  \ifdim\@tempa\p@<\p@\scalebox{\@tempa}{\usebox\pandoc@box}%
  \else\usebox{\pandoc@box}%
  \fi%
}
% Set default figure placement to htbp
\def\fps@figure{htbp}
\makeatother
\setlength{\emergencystretch}{3em} % prevent overfull lines
\providecommand{\tightlist}{%
  \setlength{\itemsep}{0pt}\setlength{\parskip}{0pt}}
\usepackage[style=apa,]{biblatex}
\addbibresource{mybib.bib}
\usepackage{ctex}

%\usepackage{xltxtra} % XeLaTeX的一些额外符号
% 设置中文字体
%\setCJKmainfont[BoldFont={黑体},ItalicFont={楷体}]{新宋体}

% 设置边距
\usepackage{geometry}
\geometry{%
  left=2.0cm, right=2.0cm, top=3.5cm, bottom=2.5cm} 

\usepackage{amsthm,mathrsfs}
\usepackage{booktabs}
\usepackage{longtable}
\makeatletter
\def\thm@space@setup{%
  \thm@preskip=8pt plus 2pt minus 4pt
  \thm@postskip=\thm@preskip
}
\makeatother
\usepackage{bookmark}
\IfFileExists{xurl.sty}{\usepackage{xurl}}{} % add URL line breaks if available
\urlstyle{same}
\hypersetup{
  pdftitle={二维多尺度问题正交分解方法及其程序实现},
  pdfauthor={李东风},
  hidelinks,
  pdfcreator={LaTeX via pandoc}}

\title{二维多尺度问题正交分解方法及其程序实现}
\author{李东风}
\date{2030年6月5日}

\begin{document}
\maketitle

{
\setcounter{tocdepth}{2}
\tableofcontents
}
\section{摘要}\label{intro}

本文围绕二维多尺度椭圆型问题的局部正交分解方法展开研究，聚焦理论构建与程序实现。此处可进一步补充研究背景、研究目标、技术路线、核心结果与结论，形成完整摘要正文。

\section{第一章 绪论}\label{chap-1}

\subsection{1.1 研究背景与意义}\label{chap-1-1}

本节用于阐述现代复合材料、多孔介质等介质在微观尺度上的非均匀性，说明其对数值模拟精度与计算代价提出的挑战，并明确本文研究的现实意义与理论价值。

\subsection{1.2 多尺度数值计算方法综述}\label{chap-1-2}

本节用于综述同质化理论、多尺度有限元法与变分多尺度法的主要思路、适用条件及局限，为后续引入LOD方法提供比较基础。

\subsection{1.3 局部正交分解（LOD）方法简介}\label{chap-1-3}

本节用于说明LOD方法的核心思想，重点强调其对周期性假设与严格尺度分离条件的弱依赖特性。

\subsection{1.4 本文主要研究内容与结构安排}\label{chap-1-4}

本节用于概述本文的研究问题、技术路线、程序实现框架以及各章节之间的逻辑关系。

\section{第二章 椭圆型多尺度问题的数学描述}\label{chap-2}

\subsection{2.1 模型问题定义}\label{chap-2-1}

本节用于给出具有强非均匀、高震荡系数的二维泊松方程模型，并写出对应弱形式与边界条件设定。

\subsection{2.2 标准有限元法的局限性}\label{chap-2-2}

本节用于展示标准有限元法在未解析尺度（Pre-asymptotic regime）下的精度失效现象，并说明引入多尺度方法的必要性。

\subsection{2.3 能量空间与有限元空间定义}\label{chap-2-3}

本节用于定义连续能量空间、粗网格空间 \(V_H\) 与细网格参考空间 \(V_h\)，为后续尺度分解与误差分析做准备。

\section{第三章 二维LOD方法的理论框架}\label{chap-3}

\subsection{3.1 尺度的分解与拟插值算子}\label{chap-3-1}

本节用于定义拟插值算子 \(I_H\) 及其关键性质，并建立解空间由粗尺度空间与细尺度核空间 \(W\) 构成的分解关系。

\subsection{\texorpdfstring{3.2 \(a\)-正交化与理想多尺度空间}{3.2 a-正交化与理想多尺度空间}}\label{chap-3-2}

本节用于构造相对于双线性型 \(a(\cdot,\cdot)\) 的正交补空间 \(V_H^{ms}\)，并给出校正算子 \(\mathcal{C}\) 的定义与作用机制。

\subsection{3.3 理想LOD方法的误差分析}\label{chap-3-3}

本节用于讨论理想LOD方法在能量模与 \(L^2\) 模下的收敛性质与最优阶结果。

\section{第四章 数值校正子的局部化技术}\label{chap-4}

\subsection{4.1 校正子的指数衰减特性}\label{chap-4-1}

本节用于说明并证明校正子在空间上的指数衰减规律，作为局部化算法有效性的理论依据。

\subsection{4.2 局部化校正算子与单元补片（Patches）}\label{chap-4-2}

本节用于构造基于单元补片 \(N_\ell(T)\) 的局部化空间 \(W_\ell(T)\)，并讨论局部化参数 \(\ell\) 在精度与计算量之间的权衡关系。

\subsection{4.3 局部化正交分解方法（LOD）的离散格式}\label{chap-4-3}

本节用于给出局部化后LOD方法的离散求解格式与实现要点。

\section{第五章 算法设计与MATLAB程序实现}\label{chap-5}

\subsection{5.1 算法总体流程与数据结构}\label{chap-5-1}

本节用于结合 \texttt{code.md} 的主脚本逻辑，说明整体算法流程、输入输出关系与核心数据结构。

\subsection{5.2 核心模块实现细节}\label{chap-5-2}

本节用于描述网格生成与加密（\texttt{refineMesh} 红化加密）、拟插值算子的矩阵表示构造方法，以及基于邻接关系的局部补片定位算法。

\subsection{5.3 校正子问题的鞍点求解格式}\label{chap-5-3}

本节用于说明局部约束变分问题的鞍点系统构造、求解策略，以及校正算子矩阵与多尺度基函数的组装过程。

\subsection{5.4 粗网格刚度矩阵的组装与求解}\label{chap-5-4}

本节用于给出粗网格层面的系统组装方法、求解流程与复杂度讨论。

\section{第六章 数值算例与应用分析}\label{chap-6}

\subsection{6.1 非周期高反差系数问题算例}\label{chap-6-1}

本节用于设置非周期随机震荡系数算例，展示方法在复杂多尺度介质下的适用性。

\subsection{6.2 收敛性验证}\label{chap-6-2}

本节用于比较不同局部化参数 \(\ell\) 下的误差表现，并与标准有限元法（FEM）进行精度与效率对比。

\subsection{6.3 结果可视化与解释}\label{chap-6-3}

本节用于展示多尺度数值解与细尺度参考解的可视化结果，并给出误差来源与物理意义分析。

\section{第七章 结论与展望}\label{chap-7}

\subsection{7.1 本文总结}\label{chap-7-1}

本节用于总结LOD方法在二维多尺度问题中的主要优势、实现结果与实验结论。

\subsection{7.2 未来研究方向}\label{chap-7-2}

本节用于讨论该方法在特征值问题、抛物线问题及随机同质化方向上的潜在扩展。

\section{附录：主要源程序清单}\label{appendix-code}

本附录用于整理本文实现所涉及的关键程序模块、函数清单与调用关系，后续可按模块补充源码片段与注释说明。

\printbibliography

\end{document}
